The circle and sheet pictures motivating finite-type AES compactification are
useful local models, but they preserve different structures.  This appendix
records the exact statements needed to keep those pictures from becoming
unproved global identifications.

\subsection{Triangular gaps and four initial branches}

\begin{proposition}[M\"obius equivalence of non-degenerate triangular gaps]
\label{prop:p0-apollonian-gap-equivalence}
Let $G$ be a specified curvilinear triangular gap bounded by three mutually
tangent circles on the Riemann sphere, with three distinct tangency points
and disjoint chosen circle interiors.  Any two such marked gaps, with matching
cyclic order, are M\"obius equivalent.  This equivalence carries every finite
generation of their standard Apollonian inscribed-circle recursion to the
corresponding generation.
\end{proposition}

\begin{proof}
Send one tangency point to infinity.  The two circles through it become
parallel lines, and the third becomes a circle tangent to both lines.
A Euclidean similarity makes their separation $2$ and the circle radius $1$;
translation along the strip fixes its centre.  Selecting the matching gap
and cyclic marking fixes the remaining choice.  This gives a common normal
form.  In that strip each of the two chosen cusp gaps has a unique circle
inside it tangent to its three bounding circles: tangency to both parallel
lines forces radius $1$ and the centre onto their midline; tangency to the
existing unit circle then places the centre at distance $2$ along that
midline, on the chosen side.  M\"obius transformations preserve
circles, incidence, tangency and containment of the chosen spherical discs,
so uniqueness identifies the inserted circle and then its three child gaps.
Induction proves the assertion for every finite generation.  The standard
recursive packing is described in \cite{GrahamEtAl2005Apollonian}.
\end{proof}

The three tangencies lie on a dual circle orthogonal to the bounding circles.
In its Poincar\'e disc, their arcs bound an ideal hyperbolic triangle; doubling
that triangle along its geodesic sides produces a three-cusp sphere.  Each
child gap has its own dual disc.  The newly inscribed circle is not thereby a
horocycle for the original disc metric, and ordinary three-cusp point
compactification does not generate an Apollonian packing.

If a central circle splits a four-sided initial gap into four triangular gaps,
subsequent standard Apollonian filling produces $4\cdot3^n$ gaps at depth
$n\ge0$.  Four initial entries do not define four-way recursion at every
stage.  A four-marked parent gap also retains marking data: the real modulus
of a marked planar quadrilateral should not be confused with the complex
cross-ratio modulus of four punctures on a sphere.

\subsection{The AES four-circle configuration is not the tangent ring}

\begin{proposition}[Incidence obstruction and an explicit deformation]
\label{prop:p0-four-circle-deformation}
The four circles of Proposition~\ref{prop:p0-puncture-tangent-cone} are not
M\"obius equivalent to four equal circles in a cyclic tangent ring with
adjacent tangencies distinct.  They are connected as circle configurations
by the continuous family
\[
 C_j(t):\quad |z-t\rho i^j|=\rho,
 \qquad j=0,1,2,3,\quad 1\le t\le\sqrt2.
\]
At $t=\sqrt2$, adding $|z|=(\sqrt2-1)\rho$ leaves four triangular gaps of
the preceding proposition.
\end{proposition}

\begin{proof}
At $t=1$ all four circles pass through $0$; adjacent ones intersect at a right
angle and opposite ones are tangent at $0$.  In the tangent ring there is no
common point, adjacent circles are tangent, and opposite circles are disjoint.
M\"obius transformations preserve concurrence and intersection angles, so
cannot identify the configurations.  At $t=\sqrt2$, adjacent centre distance
is $\sqrt2t\rho=2\rho$, while opposite centre distance is $2t\rho>2\rho$.
The distance of any centre from $0$ is $\sqrt2\rho$, equal to the sum of the
large radius and $(\sqrt2-1)\rho$; hence the new central circle is tangent
to all four.  Each adjacent pair and the central circle bound a
non-degenerate triangular gap.
\end{proof}

This is a deformation of circle incidence, not a deformation theorem for AES.
The four-circle AES of Theorem~\ref{thm:p0-four-circle-model} uses only local
arcs inside $D_r$, $r<\rho$.  Moving the entire circles does not prove that
its assignment, metric, critical set or puncture count continues along the
family.  Those data must be constructed and checked separately.

\subsection{A three-sheet model requires cuts and an extra base sheet}

Let $D$ be a disc about $O$, and write $N,W,S,E$ for its northern, western,
southern and eastern half-discs.  Every open quadrant belongs to exactly two
of these four sets.  Three sheets per quadrant therefore require an
independent copy of $D$ as well.  Assemble temporary sheets
$A=D$, $B=N\cup S$, $C=W\cup E$, with the half-disc joins specified.
Cut along the positive coordinate rays and identify banks by the
transpositions $(A\ B)$ and $(A\ C)$.  The monodromy around $O$ is a
three-cycle, with its direction depending on the declared crossing convention.

\begin{proposition}[The resulting local branched cover]
\label{prop:p0-three-sheet-cover}
The transitive monodromy just described gives a connected three-sheet cover
of $D\setminus\{O\}$.  With the pulled-back complex structure its one-point
extension has local form $\pi(w)=w^3$.  The metric completion of the pulled-back
Euclidean length metric has cone angle $6\pi$ at that point.
\end{proposition}

\begin{proof}
The fundamental group of the punctured disc is $\mathbb Z$; the action of its
generator is a three-cycle and its stabilizer subgroup is $3\mathbb Z$.
The classification of covers identifies this connected cover with the
punctured-disc power map \cite[\S1.3]{Hatcher2002}.  That map extends
holomorphically by $w=0$.  If $w=re^{i\theta}$, its metric pullback is
\[
 9r^4dr^2+9r^6d\theta^2=dR^2+9R^2d\theta^2,
 \qquad R=r^3,\quad \theta\in\mathbb R/(2\pi\mathbb Z).
\]
A circle of radius $R$ therefore has length $6\pi R$.  This is a cone length
metric; its tensor degenerates at $w=0$, even though the conformal carrier is
smooth there.
\end{proof}

Identifying only the centres of independent discs does not give this surface:
its punctured neighbourhood has several components.  If instead each
half-disc carries its own complete hyperbolic metric, $O$ lies on an
infinite-distance boundary and is not a finite seam point.  Finally the
pullback of $|z-\rho|=\rho$ is $|w^3-\rho|=\rho$, generally not a circle,
while its projection remains the original circle.  The branched cover does
not convert the AES four-circle configuration into the tangent ring.
A faithful circle-domain uniformization carrying the arithmetic data remains
an open problem.

\subsection{Other completions and the genus-two comparison}

The visual compactification $\mathbb H^2\cup S^1_\infty$ is a closed disc;
adding a direction circle by the unit tangent bundle is a different,
three-dimensional construction.  A finite-word tree has its own end space,
and identifying it with a hyperbolic visual boundary requires a specified
coding and convergence theorem.  None of these operations is the conformal
point filling of Section~\ref{sec:p0-compactification}.

A companion algebraic example makes the distinction concrete
\cite{Yuan2026ObserverMetricCompletion}.  After adjoining independent sheets
\[
 Y^2=2(E-U)(1-U^2),\qquad Z_m^2=c+U^2,
\]
the product-sign quotient $W=YZ_m$ obeys
\[
 W^2=2(E-U)(1-U^2)(c+U^2).
\]
At $E=0,c=1$, setting $w=W/\sqrt2$ yields $w^2=U^5-U$.
Its five finite branch points and infinity are distinct, so the degree-two
cover of the sphere has genus $2$ by Riemann--Hurwitz; this is the Bolza model.
The compactified object is the newly constructed algebraic curve, and the
product quotient forgets the separate sheet signs.  No canonical selection of
$c$ or faithful reconstruction of those signs follows from the equation.

Likewise, pairing all three boundary circles of two compact oriented pants
with orientation-reversing boundary maps produces a closed genus-two surface:
$\chi=-1-1=-2$.  Equality of genus supplies only a topological comparison;
matching the Bolza conformal structure, marked seams and process residues
requires further data and a separate proof.  These observations locate the
cross-programme research question without asserting an equivalence between
its distinct geometric constructions.
